\section{Costello 2013 and the Chiral Yangian Boundary}
\label{sec:costello-2013-chiral-yangian-boundary}

\providecommand{\kch}{\kappa_{\mathrm{ch}}}
\providecommand{\kcat}{\kappa_{\mathrm{cat}}}
\providecommand{\gghone}{\widehat{\mathfrak{gl}}_1}
\providecommand{\Yplus}{Y^+_{\epsilon_1,\epsilon_2}}
\providecommand{\Fgauge}{\mathcal{F}^{\mathrm{gauge}}_\Omega}

\subsection{The Four-Dimensional Theorem}

Costello's \emph{Supersymmetric gauge theory and the Yangian}
\texttt{arXiv:1303.2632} proves a perturbative theorem about the
twisted and \(\Omega\)-deformed pure \(\mathcal N=1\) gauge theory.  The
output is controlled by the Yangian \(Y(\mathfrak g)\) of the gauge Lie
algebra.  The calculation is expressed through Wilson operators and the
partition function of an integrable spin chain on a doubly-periodic
lattice.

The ambient theory is pure \(\mathcal N=1\) gauge theory on a
holomorphic-topological four-manifold of the form
\[
  \mathbb R^2_{\mathrm{top}}\times \mathbb C^{\mathrm{hol}}.
\]
The Nekrasov \(\mathcal N=2\) instanton \(\Omega\)-background and the
Maulik-Okounkov stable-envelope Yangian action are separate
constructions.  They meet the Costello construction through the same
Yangian algebra, with different geometric inputs.

\subsection{Factorisation and Vertex Algebras}

The factorisation-algebra form of the statement is compound.  Costello
2013 supplies the perturbative holomorphic-topological quantisation and
the Yangian control.  Costello-Gwilliam's factorisation-algebra theorem
identifies a holomorphic factorisation algebra on \(\mathbb C\), under
the usual smooth translation and locality hypotheses, with a vertex
algebra.  Thus the statement
\[
  \Fgauge \quad \text{is a vertex algebra on } \mathbb C
\]
uses Costello 2013 together with Costello-Gwilliam.

\subsection{Boundary Chiral Algebra Lane}

The five-dimensional or six-dimensional boundary formulation is a later
construction.  Costello-Yagi derive four-dimensional Chern-Simons theory
from an \(\Omega\)-deformed six-dimensional topological-holomorphic
theory.  Costello-Gaiotto and Costello-Dimofte-Paquette identify
boundary chiral algebras in holomorphic twists under specified boundary
conditions.  In the abelian \(\gghone\) case, the resulting boundary
vertex algebra is the affine-Yangian or \(\mathcal W_{1+\infty}\) target
after the correct double, completion, and representation-theoretic
normalisation.

The clean attribution chain is therefore:
\[
\begin{array}{p{0.30\linewidth}p{0.62\linewidth}}
Costello 2013 &
perturbative four-dimensional \(\mathcal N=1\) Yangian control;\\
Costello-Gwilliam &
holomorphic factorisation algebra to vertex algebra;\\
Costello-Yagi &
six-dimensional origin of four-dimensional Chern-Simons;\\
Costello-Gaiotto and Costello-Dimofte-Paquette &
boundary chiral algebras in holomorphic twists under stated boundary
conditions.
\end{array}
\]

\subsection{Affine Three-Space Scope}

For \(\mathbb C^3\), the cohomological Hall algebra is the positive
half
\[
  \CoHA(\mathbb C^3)\simeq Y^+(\gghone).
\]
The full affine Yangian is obtained from the Drinfeld double and Cartan
completion.  The vertex algebra \(\mathcal W_{1+\infty}\) appears after
the Fock or vacuum evaluation of that completed double.  Costello 2013
is a perturbative Yangian-control source for the gauge-theory side; it
is one input in the boundary vertex-algebra construction.

\subsection{Dimension and Completion}

The \(d=3\) CY-to-chiral output is \(E_1\)-chiral after curve
specialisation.  Braided \(E_2\) structure belongs to the Drinfeld
centre, the completed double, or a module category with its own
hypotheses.  Compact or global hCS statements require orientation and
determinant-line framing, regulator, completed local functionals, QME
solution, quartic anomaly trivialisation, and factorisation or TCFT
sewing.

\subsection{Resulting Dictionary}

The precise dictionary is:
\[
\begin{array}{p{0.30\linewidth}p{0.62\linewidth}}
twisted pure \(\mathcal N=1\) gauge theory &
Costello 2013 perturbative Yangian control;\\
factorisation algebra on \(\mathbb C\) &
Costello-Gwilliam vertex-algebra theorem;\\
five-dimensional boundary chiral algebra &
Costello-Yagi plus Costello-Gaiotto/CDP boundary construction;\\
\(\CoHA(\mathbb C^3)\) &
positive half \(Y^+(\gghone)\);\\
\(\mathcal W_{1+\infty}\) &
completed double/centre/evaluation output.
\end{array}
\]

The scalar shadows \(\kch(\mathbb C^3)=3/2\) and
\(\kcat(\mathbb C^3)=1\) belong to the CY-to-chiral and categorical
Euler lanes.  They are independent of the bibliographic attribution for
the Yangian boundary theorem.
